\section{Proof of Correctness}
\subsection{Iterative Algorithms}
\begin{itemize}
  \item \textbf{Loop Invariant}: A property that holds before each iteration of a loop.
  \item \textbf{Initialization}: The invariant holds at the start of the first iteration.
  \item \textbf{Maintenance}: If the invariant holds before an iteration, it
    holds after the iteration.
  \item \textbf{Termination}: After the loop terminates, the invariant (and the
    termination condition) implies the correctness of the algorithm.
\end{itemize}
\textcolor{ForestGreen}{e.g. Selection Sort}:
\[\left( \verb+A[1..j-1]+ \text{ is sorted} \right) \text{and }\]
\[\left( \forall x\in \verb|A[1..j-1]|,y\in \verb|A[j..n]|,\quad x\le y\right) .\] 

\subsection{Recursive Algorithms}
\begin{itemize}
  \item \textbf{Base Case}: The algorithm works for the base case(s)
    (no recursive calls).
  \item \textbf{Inductive Hypothesis}: Assume the algorithm is correct for all
    inputs of size less than $n$.
  \item \textbf{Inductive Step}: Show that the algorithm works for an input of
    size $n$, using the inductive hypothesis.
\end{itemize}
